% ===== latex-recomposed:document =====
% flm: style file flmrecomposed.f95314e6.sty
\flmUDg{x}{\BooleanTrue} \flmUDg{}{\BooleanTrue} \flmUDg{}{\BooleanFalse}and text.%%
% ===== latex-recomposed:style flmrecomposed.f95314e6.sty =====
\NeedsTeXFormat{LaTeX2e}
\ProvidesPackage{flmrecomposed.f95314e6}
% --- commands and environments the library defines for the document ---
% flm.latex.def.ifblank
% \flmIfBlankTF{<content>}{<then>}{<else>}: the branch a user definition's own
% "is this blank?" test selects — <then> where <content> holds nothing at
% all, <else> otherwise. Not a customization point: change your own
% definition, not this. It is long, so that an argument holding a paragraph
% break is tested rather than stopping the compilation, and expandable, so
% that the test never stops the expansion of a definition that has to expand.
%
% This block uses expl3, LaTeX's programming layer, which
% is part of the format and needs no package; its names read only between
% \ExplSyntaxOn and \ExplSyntaxOff. \tl_if_empty:nTF is true exactly when its
% argument holds no token at all, which is the condition FLM's own blank test
% asks: a single space is content, and a comment is gone before the argument
% exists.
\ExplSyntaxOn
\providecommand\flmIfBlankTF[3]{\tl_if_empty:nTF {#1} {#2} {#3}}
\ExplSyntaxOff
% --- user definitions ---
% \flmUDg (the document's \g): defined unless the name is already taken, so a definition of it made before this file is loaded wins; \renewcommand after loading this file changes it
\providecommand\flmUDg[2]{\IfBooleanT{#2}{V}\flmIfBlankTF{#1}{B}{N}}
\endinput
